% sec_front.tex -- title, authors, abstract, classification
\title[The $6\times6\times6$ tic-tac-toe game is a draw]{The $6\times6\times6$ tic-tac-toe game is a draw}

\author{[Author names and affiliations to be supplied]}

\subjclass[2020]{Primary 91A46; Secondary 91A24, 05C65, 05D15, 68V05}

\keywords{Tic-tac-toe, positional game, Maker--Breaker game, strong draw, pairing strategy, Hall's theorem,
computer-assisted proof}

\hypersetup{pdfkeywords={tic-tac-toe, positional game, Maker-Breaker game, strong draw, pairing strategy,
Hall's theorem, computer-assisted proof}}

\begin{abstract}
We prove that tic-tac-toe on the $6\times6\times6$ board is a draw. More precisely, in the
Maker--Breaker game on the $148$ lines of this board, Breaker, moving second, can prevent Maker from ever
occupying a whole line; by strategy stealing, the ordinary game is then a draw. The proof is
computer-assisted. A pairing strategy, the classical tool for such results, does not exist on the empty
board and cannot exist before Breaker has placed seven stones. We show that Breaker can reach, within ten
rounds, a position in which a pairing strategy takes over. Until then Breaker follows a finite table,
answers Maker's immediate threats by forced replies, and takes note of Maker's other stones only when the
table runs into them; a condition derived from Hall's theorem guarantees a pairing for every admissible
placement of the stones not yet noted. Up to the symmetries of the board, the table has $93{,}164$ entries
and ends in $106{,}723$ positions that satisfy this condition. It was found by a search and verified in
exact integer arithmetic by independently written programs. As a corollary, Breaker also wins on the
$7\times7\times7$ board, a result first obtained in a recent preprint of Cain, Raymond and
K\"allersj\"o. To our knowledge, $5\times5\times5$ is now the only $n\times n\times n$ game whose outcome
is unknown.
\end{abstract}

\maketitle

\setcounter{tocdepth}{1}
\tableofcontents
